\section{Navigation}

With an address on every dock, moving through the mesh is a matter of arithmetic, but the geometry sets hard limits on how far and how fast. A tone moves one dock per beat, no faster. After $n$ beats it has reached docks at most $n$ steps away, and nothing carries news beyond that. This is the finite light cone, and it is built into the law before any physics is read from it. The speed limit of the world is one dock per tick.

Inside hyperbolic space, distance and population pull sharply apart. The number of docks within $n$ steps grows exponentially with $n$, yet the path to the farthest of them is still only $n$ steps long. A region a short walk away already holds exponentially many docks. This is why the bulk feels vast from the inside while every point stays close to the boundary, and it is the geometric root of how much a small pattern can reach. Navigation is cheap in steps and astronomically rich in destinations.

Moving is not free, though. Carrying a pattern from one region to another costs beats, and the bath at the growing edge makes some moves one-way. A path that has shed a disturbance into the edge cannot be walked back. Travel in the mesh is therefore directed, like travel on a river, easy with the current of growth and dear against it. The same asymmetry that will become time's arrow is already felt here, as the difference between an easy journey and a hard one.
